% Independent Pell improvements; input by the satellite article.
\section{Further reductions in the Pell block}\label{sec:further-pell}

Two changes to the Pell block save four operations independently of the
digit encoding: a shorter parameter in the final Pell equation saves two,
and a pair of linear interval equations saves another two. All unknowns
remain positive integers. Signed intermediate expressions retain the
auxiliary domain of Definition~\ref{def:certificate}.

\subsection{Factoring the final Pell parameter}

Write $J=2r+1$. The remark after Lemma~2.28 of~\cite{J82} explicitly
permits replacing the parameter $a+f^2(d^2-a)$ by
\[
 G_0=a+f^2(f^2-a).
\]
This already saves one operation: if $X=(a^2-1)(ic^2)^2$, then (E16)
asserts $f^2=X+1$, and
\[
 G_0-1=(f^2-1)(f^2-a+1).
\]
The four instructions
\[
 T=f^2-(a-1),\quad G_-=XT,\quad G_+=G_-+2,
 \quad D_G=G_-G_+
\]
replace five instructions for $G_0^2-1$; both $X$ and $a-1$ are already
needed elsewhere.

This replacement preserves positive witnesses constructively. Once
$c=\psi_a(J)$ and $d=\chi_a(J)$, keep positive $f,i$ satisfying (E16).
For either the original parameter or $G_0$, one has
$G\equiv a\pmod f$, $G\equiv1\pmod c$, and $G\ge2a+1$.
Set $I=\chi_G(J)$ and $H=\psi_G(J)$. The Pell recurrence polynomials give
$I\equiv d\pmod f$, $H\equiv J\pmod c$, while growth gives $I>d$ and
$H>J$. Thus $o=(I-d)/f$ and $j=(H-J)/c$ are positive integers.
This explicitly supplies the required positive first coordinate; it does
not assume that an arbitrary signed square root in Lemma~2.28 is positive.

A signed congruence permits a further reduction. Replace (E17) by
\begin{equation}\label{eq:signed-pell-new}
 (of-d)^2=(G^2-1)(J+jc)^2+1,
 \qquad G=1+(a+1)(f^2-1).
\end{equation}
The variable $o$ is positive; only the expression $of-d$ may have either
sign.

\begin{lemma}\label{lem:signed-chi-new}
For $a>1$ and $0<k\le m$, the congruence
$\chi_a(n)\equiv\pm\chi_a(k)\pmod{\chi_a(m)}$ implies
$n\equiv\pm k\pmod{2m}$.
\end{lemma}
\begin{proof}
The positive-sign case is the usual $\chi$ step-down lemma quoted after
Lemma~2.28 of~\cite{J82}, with its stronger modulus $4m$. For the negative
sign, the Pell addition identities give
\[
 \chi_a(n+2m)\equiv-\chi_a(n)\pmod{\chi_a(m)},
\]
because $\chi_a(2m)=2\chi_a(m)^2-1$ and
$\psi_a(2m)=2\chi_a(m)\psi_a(m)$. Apply the positive-sign case at
$n+2m$ and reduce modulo $2m$.
\end{proof}

\begin{proposition}\label{prop:signed-pell-new}\label{lem:signed-pell}
Under $a>1$ and $1<J<c$ with $J$ odd, equations (E15), (E16), and
\eqref{eq:signed-pell-new} are existentially equivalent to
$c=\psi_a(J)$, with every system unknown positive.
\end{proposition}
\begin{proof}
For sufficiency, write
\[
 d=\chi_a(k),\quad c=\psi_a(k),\qquad
 f=\chi_a(m),\quad ic^2=\psi_a(m).
\]
The standard Pell divisibility fact used in the proof of Lemma~2.28,
$\psi_a(k)^2\mid\psi_a(m)\Rightarrow\psi_a(k)\mid m$, gives $c\mid m$;
in particular $0<k<c\le m$. Our parameter satisfies
\[
 G>1,\qquad G\equiv-a\pmod f,\qquad G\equiv1\pmod c.
\]
Put $I=|of-d|$ and $H=J+jc$. The last Pell equation gives a positive
index $n$ with $I=\chi_G(n)$ and $H=\psi_G(n)$. Since
$I\equiv\pm d\pmod f$ and
$\chi_{-a}(n)=(-1)^n\chi_a(n)$, Lemma~\ref{lem:signed-chi-new} gives
$n\equiv\pm k\pmod{2m}$, hence modulo $c$. On the other hand,
$G\equiv1\pmod c$ gives $H\equiv n\pmod c$, so $J\equiv\pm k\pmod c$.
The positive-sign case yields $J=k$. The other case yields $J+k=c$,
impossible because $c=\psi_a(k)\equiv k\pmod2$ and $J$ is odd.

For necessity, choose a positive Pell denominator divisible by $c^2$,
as in the construction of Lemma~2.28 of~\cite{J82}; this supplies positive
$f,i$ satisfying (E16). Put
$I=\chi_G(J)$, $H=\psi_G(J)$. Since $J$ is odd,
$I\equiv-d\pmod f$ and $H\equiv J\pmod c$. Therefore
$o=(I+d)/f>0$ and $j=(H-J)/c>0$ supply the desired witnesses.
\end{proof}

For its certificate, compute $a_-=a-1$, $a_+=a+1$, and
$A=a_-a_+$ in place of the three existing instructions jointly needed
for $a^2-1$ and $a-1$. The parameter coefficient then takes only
\[
 G_-=a_+X,\qquad G_+=G_-+2,\qquad D_G=G_-G_+.
\]
Thus \eqref{eq:signed-pell-new} saves two operations relative to the
original five-instruction parameter computation. Substituting $X=f^2-1$
uses (E16): the verifier records the exact triangular correction
\[
 F_{17}^{\rm calc}=F_{17}^{\rm source}
  +F_{16}(a+1)(G_{\rm source}+G_{\rm calc})(J+jc)^2.
\]

\subsection{Replacing the quadratic approximation by an interval}

The factor $4$ in (E10) is unnecessary for these encodings, and the
square can subsequently be removed too. The proof must establish the
evenness used in rounding before appealing to binomial divisibility.
We give that argument explicitly.

Use the notation
\[
 \begin{gathered}
 N=n,\quad R=r,\quad Y=sN^2,\quad U=wN^2,\quad M=RY,\\
 A=a=M(U+1),\quad P=p=2UM^2,\quad K=k,\quad C=c,\quad J=2R+1.
 \end{gathered}
\]
The encoding, independently of the Pell block, establishes
\[
 3<3L\le B<q\le N\le R,\qquad R,N\ge8,\qquad 0<b\le N.
\]
Moreover $B=Hb^4$ is even because admissibility fixes an even $H$.
This does not yet assert that $q$ or $N$ is even. Both $N=q^{16}$ and
the shorter encoding $N=q^8$ satisfy the argument below.

\begin{lemma}\label{lem:weak-pell-rounding-new}
Replace (E10) temporarily by
$(C-KY)^2+\eta=K^2$ with $\eta>0$, retaining the other equations.
Then all conclusions of the original Pell block still hold, and in fact
$|C/K-Y|<1/2$.
\end{lemma}
\begin{proof}
We have $U,Y\ge64$, $M\ge512$, $A\ge33280$, and
\begin{equation}\label{eq:weak-error-new}
 |C/K-Y|<1.
\end{equation}
Equation (E11) gives $K\ge R+2$, so
$C>(Y-1)K\ge63(R+2)>J$. Therefore either the source Pell construction
or Proposition~\ref{prop:signed-pell-new} gives $C=\psi_A(J)$.
By (E9), (E11), and the elementary Pell congruence,
\[
 K=\psi_P(R+1+v(P-1)),\qquad v\ge0.
\]
If $v\ge1$, Lemma~2.19 of~\cite{J82} gives
\[
 \frac CK\le
 \frac{(2A)^{2R}}{(2P-1)^{R+P-1}}
 \le\left(\frac{2A}{2P-1}\right)^{2R}<\frac12,
\]
contradicting $C/K>Y-1\ge63$. Hence $K=\psi_P(R+1)$.

Set $\xi=(U+1)^{2R}/U^R$. The elementary estimates in the proof of
Lemma~2.25 now apply to these exact Pell indices and give
\begin{equation}\label{eq:pell-ratio-new}
 \xi\left(1-\frac{R}{M(U+1)}\right)<\frac CK
 <\xi\left(1+\frac{R}{2M^2U}\right).
\end{equation}
In particular
$C/K>\xi/2>U^R/2>U^{R-1}+1$. Combining this with
\eqref{eq:weak-error-new} yields $Y>U^{R-1}$ and therefore
\begin{equation}\label{eq:a-growth-new}
 A=RY(U+1)>RU^R.
\end{equation}

Put $W=bw$. We have
$2^{3J}\le8N^{2R}\le RU^R<A$ and $W^3\le U^R<A$.
Equation (E14), with $C=\psi_A(J)$ and $d=\chi_A(J)$, is exactly the
$\chi$ congruence of Lemma~2.22. Thus $W=2^J$, which implies that
$b$ is a power of two and
\begin{equation}\label{eq:u-growth-new}
 U=N^2w\ge Nbw=NW\ge8\cdot2^J>4\cdot2^{2R}.
\end{equation}
Next, $\kappa<C$, (E19), and (E20), together with $0<L<J<A$, give
$\kappa=\psi_A(L)$ by Lemma~2.23. By \eqref{eq:a-growth-new},
\[
 B^{3L}\le B^B\le N^N\le U^R<A,
 \qquad q^3\le N^N\le U^R<A.
\]
Equation (E18) and the $\chi$ form of Lemma~2.22 therefore give $q=B^L$.
Since $B$ is even, $q$ and $N$ are now even. These two exponential
arguments precede the rounding step; they do not invoke Lemma~2.26
with a weakened hypothesis.

From $\xi/2<C/K<Y+1$ we get $\xi<2Y+2<3Y$. Equation
\eqref{eq:pell-ratio-new} therefore gives
\[
 \left|\frac CK-\xi\right|
 <\frac{R\xi}{M(U+1)}
 =\frac{\xi}{Y(U+1)}<\frac3{65}<\frac14.
\]
The binomial expansion of Lemma~2.24 and \eqref{eq:u-growth-new} give
\[
 0\le\xi-\lfloor\xi\rfloor<\frac14,
 \qquad \lfloor\xi\rfloor\equiv\binom{2R}{R}\pmod U.
\]
Hence $|Y-\lfloor\xi\rfloor|<3/2<2$. Both integers are even:
$Y=sN^2$, $U=wN^2$, and
$\binom{2R}{R}=2\binom{2R-1}{R-1}$. Thus $Y=\lfloor\xi\rfloor$.
This proves $N^2\mid\binom{2R}{R}$, and the last two error estimates
also give $|C/K-Y|<1/2$, restoring the original stronger bound.
\end{proof}

\begin{proposition}\label{prop:linear-pell-interval-new}\label{lem:linear-ratio}
In the admissible encodings considered here, (E10) may be replaced by
\begin{equation}\label{eq:linear-pell-interval-new}
 c=ksn^2+\eta,\qquad k=\eta+\zeta,
 \qquad \eta,\zeta>0.
\end{equation}
This replacement saves two operations and adds one positive unknown
and one free equality test.
\end{proposition}
\begin{proof}
The two equations give $0<C/K-Y<1$, so Lemma~\ref{lem:weak-pell-rounding-new}
proves sufficiency. For necessity, an original solution satisfies the
weaker bound too: $K^2-(C-KY)^2$ is a positive integer and supplies its
slack. Applying that lemma gives
$Y=\lfloor\xi\rfloor$ and $\xi<Y+1/4<2Y$. The lower estimate in
\eqref{eq:pell-ratio-new} implies
\[
 \xi-\frac CK<\frac{\xi}{Y(U+1)}<\frac2{U+1}<\frac2U.
\]
The positive term of order $1/U$ in the binomial expansion gives
\[
 \xi-Y\ge\frac{\binom{2R}{R-1}}{U}
 \ge\frac{2R}{U}>\frac2U.
\]
Consequently $C/K>Y$. Together with $|C/K-Y|<1/2$, this proves that
$\eta=C-KY$ and $\zeta=K-\eta$ are positive integers.
The certificate computes $ksn^2$, adds $\eta$, and computes
$\eta+\zeta$: three operations in place of the previous five.
\end{proof}

The proofs above preserve the other positive witnesses; $j,o$ may be
re-chosen for the new Pell parameter, and $\eta,\zeta$ supply the new
interval. The residual verifier checks all changed equations and the
two triangular identities exactly. Positivity and the equivalence of
the encodings depend on the mathematical arguments, not on polynomial
expansion alone.
